% from ega4-11-fr.tex
\oldpage[IV-3]{135}

\begin{proposition}[11.3.7]
\label{IV.11.3.7-fr}
Soient $f:X\to Y$ un morphisme localement de présentation finie,
$\mathcal F$, $\mathcal F'$ deux $\mathcal O_X$-Modules quasi-cohérents de
présentation finie, $u:\mathcal F'\to\mathcal F$ un
$\mathcal O_X$-homomorphisme, $x$ un point de $X$ et $y=f(x)$~; on suppose que
$\mathcal F$ est $f$-plat au point $x$. Les conditions suivantes sont alors
équivalentes~:
\begin{enumerate}[label=\emph{\alph*)}]
\item On a $(\ker u)_x=0$ et $\operatorname{Coker}u$ est un
$\mathcal O_X$-Module $f$-plat au point $x$.
\item L'homomorphisme
\[
(u\otimes1)_x:(\mathcal F'\otimes_{\mathcal O_Y}k(y))_x
\longrightarrow(\mathcal F\otimes_{\mathcal O_Y}k(y))_x
\]
est injectif.
\end{enumerate}
Si de plus $\mathcal F$ est $f$-plat, l'ensemble des points $x$ vérifiant les
conditions équivalentes précédentes est ouvert dans $X$.
\end{proposition}

\begin{proof}
La condition a) entraîne b) en vertu de $(\mathbf 0_{\mathrm I},6.1.2)$, sans
hypothèse sur $\mathcal F$. Pour démontrer la réciproque, on peut se borner au
cas où $X$ et $Y$ sont affines, puis, en vertu de (8.9.1) et (11.2.7), supposer
que l'on a
\[
X=X_0\times_{Y_0}Y,\quad f=f_0\times1,\quad
\mathcal F=\mathcal F_0\otimes_{\mathcal O_{Y_0}}\mathcal O_Y,
\]
\[
\mathcal F'=\mathcal F'_0\otimes_{\mathcal O_{Y_0}}\mathcal O_Y,
\qquad u=u_0\otimes1_Y,
\]
où $Y_0$ est noethérien, $f_0:X_0\to Y_0$ un morphisme de type fini,
$\mathcal F_0$, $\mathcal F'_0$ deux $\mathcal O_{X_0}$-Modules cohérents,
$u_0:\mathcal F'_0\to\mathcal F_0$ un homomorphisme~; en outre, si $x_0$ est la
projection de $x$ dans $X_0$, on peut supposer que $\mathcal F_0$ est
$f_0$-plat au point $x_0$. Posons $y_0=f_0(x_0)$~; en vertu de la transitivité
des fibres $(\mathbf I,3.6.4)$, le morphisme projection
$f^{-1}(y)\to f_0^{-1}(y_0)$ est fidèlement plat (2.2.13), et comme
\[
(u\otimes1)_x=((u_0\otimes1)_{x_0})\otimes1_{k(y)},
\]
l'hypothèse que $(u\otimes1)_x$ est injectif entraîne qu'il en est de même de
$(u_0\otimes1)_{x_0}$ (2.2.7). Or, cela entraîne, par
$(\mathbf 0_{\mathrm{III}},10.2.4)$ appliqué aux anneaux locaux noethériens
$\mathcal O_{Y_0,y_0}$ et $\mathcal O_{X_0,x_0}$ et aux
$\mathcal O_{X_0,x_0}$-modules $(\mathcal F'_0)_{x_0}$ et
$(\mathcal F_0)_{x_0}$ (dont le second est plat sur $\mathcal O_{Y_0,y_0}$),
que l'on a $(\ker u_0)_{x_0}=0$ et que $\operatorname{Coker}u_0$ est
$f_0$-plat au point $x_0$. On déduit d'abord de là que
$\operatorname{Coker}u$ est $f$-plat au point $x$ (2.1.4)~; en vertu de
l'exactitude à droite du produit tensoriel, on a d'ailleurs
\[
\operatorname{Coker}u=(\operatorname{Coker}u_0)
 \otimes_{\mathcal O_{Y_0}}\mathcal O_Y~;
\]
appliquant alors $(\mathbf 0_{\mathrm I},6.1.2)$ à la suite (exacte par
hypothèse)
\[
0\longrightarrow(\mathcal F'_0)_{x_0}\longrightarrow
(\mathcal F_0)_{x_0}\longrightarrow(\operatorname{Coker}u_0)_{x_0}
\longrightarrow0,
\]
on en déduit que $u_x$ est un homomorphisme injectif.

Enfin, il résulte de (11.1.2) que l'ensemble $U_0$ des points $z_0\in X_0$ tels
que le morphisme $(u_0\otimes1)_{z_0}$ soit injectif est ouvert~; par platitude
on en déduit que pour tout $z\in X$ au-dessus de $z_0\in U_0$ le morphisme
$(u\otimes1)_z$ est injectif, donc l'ensemble de ces points contient l'image
réciproque $U$ de $U_0$ dans $X$ et est par suite un voisinage de $x$, ce qui
achève la démonstration.
\end{proof}

\begin{theorem}[11.3.8]
\label{IV.11.3.8-fr}
Soient $f:X\to Y$ un morphisme localement de présentation finie,
$\mathcal F$ un $\mathcal O_X$-Module quasi-cohérent de présentation finie,
$(g_i)_{1\leq i\leq n}$ une suite de sections de $\mathcal O_X$ au-dessus de
$X$~; posons
\[
\mathcal F_i=\mathcal F\big/\sum_{j=1}^{i}g_j\mathcal F
\quad\text{pour }0\leq i\leq n
\quad(\text{avec }\mathcal F_0=\mathcal F).
\]
Soient $x$ un point de $X$, $y=f(x)$, et posons
$X_y=f^{-1}(y)=X\times_Y\operatorname{Spec}(k(y))$,
$\mathcal F_y=\mathcal F\otimes_{\mathcal O_Y}k(y)$, qui est un
$\mathcal O_{X_y}$-Module. On suppose que les $(g_i)_x$ appartiennent à l'idéal
maximal $\mathfrak m_x$. Les conditions suivantes sont équivalentes~:
\begin{enumerate}[label=\emph{\alph*)}]
\item La suite $((g_i)_x)$ est $\mathcal F_x$-régulière et les $\mathcal F_i$
$(0\leq i\leq n)$ sont $f$-plats au point $x$.
\item La suite $((g_i)_x)$ est $\mathcal F_x$-régulière et $\mathcal F_n$ est
$f$-plat au point $x$.
\item[\emph{b$'$)}] Il existe un voisinage $U$ de $x$ tel que la suite
$(g_i|U)$ soit $(\mathcal F|U)$-quasi-régulière, et $\mathcal F_n$ est
$f$-plat au point $x$.

